% from ega4-14-fr.tex
\oldpage[IV-3]{200}
localement noethériens et $Y$ irréductible) admettent «~suffisamment~» de
\emph{quasi-sections}, i.e. au voisinage d'un point fermé $x$ d'une fibre
$f^{-1}(y)$, il y a un sous-préschéma fermé $X'$ de $X$ contenant $x$ et tel que
la restriction $X'\to Y$ de $f$ soit un morphisme \emph{quasi-fini} (donc à fibres
\emph{discrètes}) et dominant (14.5.3).

Au \S~15 on étudie diverses propriétés des fibres des morphismes universellement
ouverts $f:X\to Y$, localement de type fini, notamment lorsque $X$ et $Y$ sont
localement noethériens. On obtient ainsi en particulier un critère pour qu'un
point $x\in X$ n'appartienne qu'à \emph{une seule} composante irréductible de $X$,
en termes de propriétés de \emph{la fibre} $f^{-1}(f(x))$ de $x$~: il suffit que
$Y$ soit géométriquement unibranche (par exemple normal) au point $f(x)$ et que
$f^{-1}(f(x))$ soit géométriquement ponctuellement intègre au point $x$ (15.3.3)~;
si de plus $Y$ est localement intègre au point $f(x)$, $f$ est \emph{plat} au point
$x$ et $X$ localement intègre au point $x$. On étudie aussi la variation du
\emph{nombre géométrique de composantes connexes} d'une fibre $f^{-1}(y)$~; par
exemple, si $f$ est universellement ouvert et \emph{propre}, et les fibres de $f$
géométriquement réduites, ce nombre est \emph{localement constant} (15.5.7).
Enfin, lorsque $f$ admet une \emph{section} $g$ (ce qui sera le cas lorsque $X$ est
un $Y$-\emph{schéma en groupes}) et que pour tout $y\in Y$, on désigne par $X_y^0$
la composante connexe de la fibre $X_y=f^{-1}(y)$ au point $g(y)$ («~composante
neutre~» dans le cas des groupes), on étudie la réunion $X^0$ des $X_y^0$ pour
$y\in Y$, et on montre (15.6.4) que si $f$ est universellement ouvert et les fibres
$X_y$ géométriquement réduites, alors $X^0$ est un ensemble \emph{ouvert} dans $X$.

\subsection{Morphismes ouverts}
\label{subsection:IV.14.1-fr}

\begin{env}[14.1.1]
\label{IV.14.1.1-fr}
Rappelons (1.10.2) qu'une application continue $\psi:X\to Y$ est dite
\emph{ouverte en un point} $x\in X$ si l'image par $\psi$ de \emph{tout} voisinage
de $x$ dans $X$ est un voisinage de $\psi(x)$ dans $Y$.

On notera que cela \emph{n'implique pas} qu'il existe un système fondamental de
voisinages de $x$ dont les images soient \emph{ouvertes} dans $Y$.
\end{env}

\begin{proposition}[14.1.2]
\label{IV.14.1.2-fr}
Soient $X$, $Y$ deux espaces topologiques, $\psi:X\to Y$ une application
continue, $x$ un point de $X$, $y=\psi(x)$.
\begin{enumerate}[label=(\roman*)]
  \item Si $\psi$ est ouverte en $x$, alors pour toute partie $Y'$ de $Y$ contenant
  $\psi(x)$ la restriction $\psi^{-1}(Y')\to Y'$ de $\psi$ à $Y'$ est ouverte au
  point $x$.
  \item Supposons que $Y$ soit réunion d'une famille localement finie de parties
  fermées $(Y_i)$ et que pour tout $i$ tel que $\psi(x)\in Y_i$, la restriction
  $\psi^{-1}(Y_i)\to Y_i$ de $\psi$ soit ouverte au point $x$~; alors $\psi$ est
  ouverte au point $x$.
  \item Soient $\gamma:X'\to X$ une application continue, $x'$ un point de $X'$~;
  si l'application composée $\psi\circ\gamma:X'\to Y$ est ouverte au point $x'$,
  $\psi$ est ouverte au point $\gamma(x')$.
\end{enumerate}
\end{proposition}

\begin{proof}
Si $U$ est un voisinage de $x$, on a
$\psi(U\cap\psi^{-1}(Y'))=\psi(U)\cap Y'$~; d'où aussitôt (i). Pour prouver (ii),
notons qu'il y a un voisinage $W$ de $\psi(x)$ dans $Y$ ne rencontrant que les
parties fermées $Y_i$ en nombre \emph{fini} qui \emph{contiennent} $\psi(x)$~; donc
$W$ est réunion des $W\cap Y_i$ pour ces indices. Or, si $U$ est un voisinage de
$x$ tel que $\psi(U)\cap Y_i$ soit un voisinage de $\psi(x)$ dans $Y_i$, il existe
un voisinage $V\subset W$ de $\psi(x)$ dans $Y$ tel que pour tous les $i$ tels que
$\psi(x)\in Y_i$,
